% ==================== PARTIE E — SYNTHÈSE ====================
% Étape de conception : intégration (fonctions + boucle + décisions) et recul critique.
%
%   E1  compléter un programme utilisant determiner_mode (prototype, appel, codes)
%   E2  écrire tarif_applicable (réutilisation possible de determiner_mode)
%   E3  critique de conception

\element{synthese-utilisation}{
    \begin{question}{synthese-utilisation}
        On dispose de la fonction
        \lstinline[language=C]|int determiner_mode(float puissance_kw, int cable_verrouille)|,
        que l'on considère \textbf{correcte}. Elle retourne \lstinline[language=C]|-1| si le
        câble n'est pas verrouillé, \lstinline[language=C]|0| en mode lent,
        \lstinline[language=C]|1| en mode rapide, \lstinline[language=C]|2| en mode
        ultra-rapide et \lstinline[language=C]|3| si la demande dépasse 150 kW. Sa définition
        est placée \textbf{après} le \lstinline[language=C]|main| ; son prototype est rappelé
        en tête du fichier.

        Compléter les cadres pointillés pour que le programme traite les
        \lstinline[language=C]|nb_sessions| demandes de la journée, compte les sessions
        refusées (câble non verrouillé ou demande hors domaine) et affiche le mode des autres.
        Le choix selon le code retourné est réalisé avec une structure
        \lstinline[language=C]|switch|.

        \lstinputlisting[language=c,escapeinside={(*}{*)}]{sources/codes/e1-utilisation-mode-a-completer.c}

        \evaluationProf[3]
    \end{question}
}

\element{synthese-evolution}{
    \begin{question}{synthese-evolution}
        L'exploitant met en place une tarification heures creuses : de 22 h à 5 h inclus, le
        tarif est de 0,15 €/kWh ; le reste du temps, il vaut 0,22 €/kWh jusqu'à 50 kW inclus
        et 0,30 €/kWh au-delà de 50 kW.

        Compléter le corps de la fonction \lstinline[language=C]|tarif_applicable| qui retourne
        le tarif en €/kWh correspondant à une heure de départ de session (entier de 0 à 23) et
        à une puissance donnée. Elle peut utiliser les fonctions déjà écrites dans le sujet,
        par exemple \lstinline[language=C]|determiner_mode| (codes rappelés à la question
        précédente). Elle ne doit rien afficher.

        \lstinputlisting[language=c]{sources/codes/e2-tarif-debut.c}
        \vspace*{-.5em}
        \evaluationProfGrille{4}{10}
        \lstinputlisting[language=c]{sources/codes/d4-fin.c}
    \end{question}
}

\element{synthese-critique}{
    \begin{question}{synthese-critique}
        En reprenant le programme de traitement des sessions de cette partie (utilisation de
        \lstinline[language=C]|determiner_mode|), identifier et expliquer
        \textbf{deux} faiblesses de conception (lisibilité, valeurs numériques écrites en dur,
        robustesse face à des saisies inattendues, absence de vérification, découpage en
        fonctions, comparaison de nombres à virgule\ldots), en indiquant pour chacune comment
        y remédier.

        \evaluationProfGrille{2}{6}
    \end{question}
}
